% Generated draft: requires mathematical and source review.
\documentclass[11pt]{article}
\usepackage[margin=25mm]{geometry}
\usepackage{fontspec}
\setmainfont{Latin Modern Roman}
\usepackage{amsmath,amssymb,mathtools,unicode-math}
\setmathfont{Latin Modern Math}
\usepackage{enumitem,needspace,xurl,hyperref}
\setlength{\parindent}{0pt}
\setlength{\parskip}{5pt}
\setlength{\emergencystretch}{3em}
\newcommand{\N}{\mathbb{N}}
\newcommand{\Z}{\mathbb{Z}}
\newcommand{\Q}{\mathbb{Q}}
\newcommand{\R}{\mathbb{R}}
\newcommand{\C}{\mathbb{C}}
\newcommand{\F}{\mathbb{F}}
\newcommand{\A}{\mathbb{A}}
\newcommand{\PP}{\mathbb{P}}
\newcommand{\E}{\mathbb{E}}
\newcommand{\eps}{\varepsilon}
\newcommand{\dd}{\,\mathrm{d}}
\begin{document}
\section*{UnsolvedMath ID: 30001092}
\textbf{Source problem code:} catalogue entry
\subsection{Definitions and mathematical statement}
Let \(\Gamma\) be a countable group acting by probability-measure-preserving transformations on a standard probability space \((X,\mu)\). Its orbit equivalence relation is
\[
 R_\alpha=\{(x,\alpha(\gamma)x):x\in X,\ \gamma\in\Gamma\}.
\]
All relations and properties are understood modulo null sets. The action is essentially free when \(\mu(\{x:\alpha(\gamma)x=x\})=0\) for every \(\gamma\ne1\).

A measured free-product decomposition of a countable p.m.p. equivalence relation expresses it as generated by subrelations \(R_i\) with no nontrivial reduced alternating cycles among the factors. Freely indecomposable (FI) means that every such decomposition is trivial in the source's measured sense. That sense must account for the allowed inessential decompositions or restrictions, and is not defined in the extraction. It cannot safely be replaced merely by a group-theoretic free-product condition on \(\Gamma\), or by an unqualified prohibition of every syntactic decomposition into subrelations.

The question asks whether there is one countable group and two actions, possibly on different probability spaces,
\[
 \Gamma\curvearrowright^{\alpha}(X,\mu),\qquad
 \Gamma\curvearrowright^{\beta}(Y,\nu),
\]
for which
\[
 R_\alpha\text{ is FI},\qquad R_\beta\text{ is not FI}.
\]
The group is the same; the actions and their orbit relations may differ. Non-FI requires a genuinely nontrivial measured free-product decomposition in the intended convention, not simply nonisomorphism to the first relation.

The original and clean statements say actions without explicitly requiring essential freeness. The supplied literature assessment instead explicitly asks for free p.m.p. actions. Thus there are two different quantified readings: the literal existence problem above, and the intended stronger version requiring both \(\alpha\) and \(\beta\) to be essentially free. This distinction is retained rather than silently resolved. Ergodicity and nonatomicity are also not explicitly supplied.

In the usual terminology invoked by the assessment, a group is measurably freely indecomposable (MFI) when every action in the intended free p.m.p. class has FI orbit relation. The free-action version here asks for a group with at least one FI action and at least one non-FI action, hence not MFI; it is not asking only whether some group fails MFI. The assessment dated 2026-08-21 reports FI criteria and MFI results as partial information, with no general example or impossibility theorem located for the intended question.
\subsection{Short English statement}
Can one countable group have two p.m.p. actions, one with freely indecomposable orbit relation and one without? The assessment intends both actions to be free, a hypothesis absent from the original wording; the measured definition of a trivial free-product decomposition also needs the full source.
\subsection{Sources}
\begin{itemize}
\item Local UnsolvedMath registry, permanent ID \(30001092\), exact source code \texttt{catalogue entry}, proposed year 2008; no individual proposer is supplied.
\item \emph{Von Neumann Algebras and Ergodic Theory of Group Actions (2009)}, Oberwolfach Reports, DOI 10.4171/owr/2008/49, pp. 2763--2814. \url{https://doi.org/10.4171/owr/2008/49}.
\item Alvarez--Gaboriau, ``free products and orbit equivalence,'' \url{https://arxiv.org/abs/0806.2788}.
\item The supplied statement status is \texttt{exact}; its verification says the internally readable statement was retained byte-for-byte, rather than asserting a source-level recovery of missing hypotheses. The open label and assessment checked 2026-08-21 are supplied dated curation, not independently verified; its explicitly free-action scope requires reconciliation with the literal text.
\end{itemize}
\end{document}
